
\chapter{Conclusão}
\label{cap:conclusao}

Este trabalho investigou a automação da elicitação de requisitos de software a partir do áudio de reuniões, propondo e implementando um sistema multiagente baseado em inteligência artificial generativa. Partindo da constatação de que a elicitação manual é cara, trabalhosa e suscetível à perda de informação, o sistema decompõe o problema em quatro agentes especializados que, em conjunto, transformam uma gravação de reunião em documentação de requisitos estruturada, com dois formatos de saída selecionáveis (o padrão IEEE~830 e um formato específico de cliente). A arquitetura foi validada em um caso real, por meio de uma parceria com um órgão público, e avaliada de forma quantitativa em cada um de seus estágios.

\section{Contribuições}

Em retrospecto, este trabalho entregou as seguintes contribuições:

\begin{enumerate}
    \item[1.] \textbf{Um sistema multiagente completo e de código aberto} para a elicitação de requisitos de ponta a ponta a partir do áudio de reuniões, integrando transcrição automática, identificação de personas, geração de histórias de usuário e produção do documento de requisitos em dois formatos.
    \item[2.] \textbf{Uma validação no mundo real}, sobre um projeto de software de um órgão público, utilizando como referência um documento de requisitos oficial produzido manualmente, cenário raramente disponível na literatura, que tende a se apoiar em dados sintéticos.
    \item[3.] \textbf{Um arcabouço de avaliação quantitativa} que cobre a acurácia da transcrição, a qualidade da extração de informação, a qualidade percebida do documento por profissionais e a cobertura frente ao documento oficial, substituindo a avaliação qualitativa preliminar.
    \item[4.] \textbf{Uma análise empírica da propagação de erro entre agentes}, tratando como resultado de primeira ordem o modo como erros de um estágio cascateiam para os artefatos seguintes, acompanhada da discussão sobre o compromisso entre custo e qualidade na seleção de modelo por agente.
\end{enumerate}

% TODO: ao fechar os experimentos, acrescentar aqui um parágrafo sintetizando
% os principais resultados numéricos (WER/CER, F1, alfa de Cronbach, cobertura)
% e o que eles indicam sobre a viabilidade do sistema.

\section{Limitações}

Os resultados deste trabalho devem ser interpretados à luz de algumas limitações. A validação no mundo real apoia-se em um \textbf{único caso real}, o que restringe a generalização das conclusões para outros domínios e organizações. A avaliação da qualidade percebida do documento depende de uma \textbf{amostra reduzida} de profissionais, o que limita a força das inferências estatísticas e exige cautela na interpretação. Observa-se, ainda, que o sistema tende a gerar \textbf{menos itens} (em especial regras de negócio e casos de uso) do que o documento produzido manualmente, refletindo a dificuldade de capturar regras frequentemente ditas de forma implícita nas reuniões.

Há também limitações de natureza técnica. Erros de transcrição \textbf{propagam-se} para os artefatos a jusante, podendo comprometer o documento final, um modo de falha inerente a \textit{pipelines} sequenciais. O sistema depende de \textbf{APIs comerciais} de transcrição e de modelos de linguagem, o que impõe custos e questões de disponibilidade e reprodutibilidade. Por fim, a avaliação concentra-se no \textbf{português} e em um domínio específico (planejamento orçamentário governamental), de modo que o desempenho em outros idiomas e domínios permanece em aberto.

\section{Trabalhos Futuros}

A partir das limitações identificadas, vislumbram-se as seguintes direções para trabalhos futuros:

\begin{itemize}
    \item \textbf{Ampliar a validação}, processando um número maior de reuniões do \textit{corpus} AMI e, idealmente, casos reais adicionais de outros domínios, para fortalecer a validade externa.
    \item \textbf{Mitigar a propagação de erro}, por exemplo, com tratamento dedicado de vocabulário técnico e de siglas de projeto, e com etapas de pós-edição ou de verificação automática entre os agentes.
    \item \textbf{Aumentar a cobertura} de regras de negócio e de casos de uso, investigando estratégias de \textit{prompting} ou de pós-processamento que recuperem requisitos expressos de forma implícita.
    \item \textbf{Avaliar a robustez} do agente de identificação em transcrições sem rótulos de falante, medindo a dependência do sistema em relação à qualidade da segmentação por interlocutor.
    \item \textbf{Explorar diferentes modelos de linguagem por agente}, quantificando de forma sistemática o compromisso entre custo, latência e qualidade.
    \item \textbf{Incorporar revisão humana no laço} (\textit{human-in-the-loop}), permitindo que o analista valide e ajuste os artefatos intermediários antes da geração do documento final.
\end{itemize}

\section{Disponibilidade de Artefatos}
O código-fonte, os \textit{scripts} de avaliação e os artefatos gerados nesta pesquisa estão publicamente disponíveis em \url{https://github.com/gabrielsfg/PFC1}. % TODO: definir uma licença para o repositório (recomendação: MIT).
